\item \model{Magistral}

\paragraph*{جایگاه مدل‌های پایه و هندسه فضای وزن‌ها}
مدل‌های \model{Magistral} \cite{magistral_2025} (شامل نسخه \en{Small} با ۲۴ میلیارد پارامتر و \en{Medium}) نشان می‌دهند که چگونه قابلیت‌های بنیادین زبانی، استدلال و درک چندوجهی در مرحله پیش‌آموزش می‌توانند بستری پایدار برای رفتارهای پیشرفته فراهم آورند. مدل‌های پایه مورد استفاده در این پروژه، نسخه‌های پیش‌آموزش‌دیده متراکم \model{Mistral Small 3} و \model{Mistral Medium 3} هستند.

\paragraph*{تحلیل مؤلفه‌های اصلی (\en{PCA}) در زیرفضای وزن‌های پیش‌آموزش}
برای درک چگونگی رفتار بازنمایی‌های پیش‌آموزش در مواجهه با فرآیندهای بهینه‌سازی استدلالی، تغییرات بردار وزن‌های مدل در طول بهینه‌سازی تحلیل شده است. ماتریس وزن‌های چک‌پوینت‌های متوالی به صورت $X \in \mathbb{R}^{T \times W}$ گردآوری گردید ($T$ تعداد گام‌های ذخیره‌شده و $W$ تعداد کل وزن‌های شبکه). با اعمال الگوریتم لانچوس-آرنولدی (\en{Lanczos-Arnoldi Iteration})، دو بردار ویژه اول ماتریس کوواریانس $X^T X$ محاسبه و بهنجار شدند ($c_1, c_2$):
\begin{equation}
    w(\alpha_1, \alpha_2) = w^* + \alpha_1 c_1 + \alpha_2 c_2
\end{equation}
با سنجش تغییرات برازش بر روی این ابرصفحه دو بعدی، مشاهده شد که یادگیری استدلال یک خط سیر پیوسته و جهت‌دار در امتداد یک بُعد به نام «طول استدلال» را طی می‌کند. میزان پاداش خام به دست آمده ارتباط مستقیمی با طول توالی تولیدی نشان می‌دهد:
\begin{equation}
    R(y) = a \cdot \ln(\operatorname{Length}(y)) + b
\end{equation}
این تحلیل اثبات می‌کند که ساختار هندسی فضای وزن‌های حاصل از پیش‌آموزش به هیچ عنوان دچار تخریب یا تغییرات فاز ناگهانی نمی‌شود، بلکه قابلیت‌های مدل در راستای بردارهای ویژه کم‌بعد فعال می‌گردد.

\paragraph*{ماندگاری توانایی‌های چندوجهی (\en{Multimodal Free Lunch})}
مدل‌های پایه با انکودرهای تصویری پیش‌آموزش‌دیده تلفیق شده‌اند. ارزیابی‌ها نشان داد که حتی اگر آموزش‌های بعدی صرفاً بر روی دادگان متنی متمرکز باشند، ساختار بازنمایی چندوجهی نه تنها دست‌نخورده باقی می‌ماند، بلکه به دلیل عمق تفکر استنتاجی، توانایی استدلال روی تصاویر در بنچمارک‌هایی نظیر \en{MMMU-Pro-Vision} تا ۱۲ درصد بهبود می‌یابد.